\chapter{PENDAHULUAN}

\section{Latar Belakang Penelitian}
Penyakit merupakan salah satu faktor utama yang memengaruhi produktivitas dan kualitas tanaman. Secara garis besar, penyakit tanaman dapat diklasifikasikan menjadi dua kategori, yakni penyakit abiotik dan penyakit biotik \parencite{bhargavaPlantLeafDisease2024}. Penyakit abiotik disebabkan oleh faktor lingkungan non-hidup, seperti kondisi tanah, suhu, atau kelembaban yang tidak ideal, yang dapat melemahkan tanaman dan membuatnya lebih rentan. Sementara itu, penyakit biotik disebabkan oleh organisme hidup penyebab penyakit (patogen), seperti jamur, bakteri, atau virus. Penyakit abiotik umumnya lebih responsif terhadap perbaikan kondisi lingkungan seperti penyesuaian penyiraman dan pemberian nutrisi, sedangkan penyakit biotik lebih sulit ditangani karena membutuhkan identifikasi yang akurat terhadap patogen penyebabnya \parencite{bhargavaPlantLeafDisease2024}. 

\textcite{fao2020climate} melaporkan penyakit tanaman menjadi salah satu ancaman utama kerugian hasil panen, di mana sekitar \qtyrange{20}{40}{\percent} hasil panen global hilang setiap tahunnya akibat serangan hama dan penyakit tanaman. \textcite{gullinoPositioningPlantHealth2024} menyebutkan bahwa tekanan lingkungan, khususnya perubahan iklim, semakin memperparah ancaman ini. Hal ini karena pemanasan global dan perubahan pola cuaca mengubah siklus hidup patogen serta sebaran geografisnya, menyebarkan penyakit ke wilayah-wilayah baru, dan memperbesar risiko wabah penyakit. Kondisi ini menunjukkan bahwa penyakit tanaman merupakan ancaman serius bagi produktivitas pertanian, yang menuntut perhatian dan upaya pengendalian. Upaya pengendalian ini diawali dengan identifikasi yang tepat terhadap gejala yang muncul pada tanaman.

Patogen penyebab penyakit tanaman biasanya menyerang daun terlebih dahulu sebelum menginfeksi seluruh tanaman \parencite{sharmaClassificationPlantLeaf2020}. Gejala penyakit tanaman yang paling mudah dikenali oleh mata manusia umumnya berupa bercak pada daun, meskipun ada penyakit yang tidak menampakkan gejala pada daun sama sekali dan ada pula yang baru terlihat pada tahap lanjut setelah tanaman mengalami kerusakan signifikan \parencite{upadhyayDeepLearningComputer2025}. Dalam praktik tradisional berarti identifikasi penyakit dilakukan dengan mengamati gejala tersebut secara visual oleh petani atau ahli yang memakan waktu dan bergantung pada ketersediaan tenaga ahli \parencite{bhargavaPlantLeafDisease2024}. Teknik laboratorium dapat memastikan identifikasi patogen penyebab penyakit secara spesifik, namun membutuhkan biaya yang relatif mahal dan memerlukan peralatan khusus serta keahlian tinggi \parencite[Sahu dkk., 2021 dalam ][]{upadhyayDeepLearningComputer2025}. Keterbatasan metode identifikasi tradisional dan laboratorium mendorong kebutuhan akan pendekatan baru yang lebih cepat, akurat, dan dapat diakses secara luas dalam mendeteksi penyakit tanaman.

Terobosan dalam \textit{artificial intelligence} (AI) kini mendorong identifikasi penyakit tanaman melalui pendekatan \textit{deep learning} dan \textit{computer vision} \parencite{upadhyayDeepLearningComputer2025}. \textit{Computer vision} adalah bidang AI yang memungkinkan komputer memahami informasi dari data visual seperti gambar atau video. Dalam identifikasi penyakit tanaman, teknologi ini dimanfaatkan untuk mengenali gejala visual melalui ciri seperti perubahan warna, bentuk area gejala, dan tekstur guna membedakan tanaman sehat dan terinfeksi. Seiring perkembangannya, \textit{computer vision} sering dikombinasikan dengan \textit{deep learning} \parencite{bhargavaPlantLeafDisease2024, upadhyayDeepLearningComputer2025}. Daun merupakan bagian tanaman yang paling memperlihatkan gejala penyakit, sehingga sebagian besar penelitian memanfaatkan gambar daun sebagai data pelatihan. Model \textit{deep learning} dilatih untuk mengenali pola visual, termasuk variasi warna, bentuk, dan tekstur daun yang menjadi indikator penyakit \parencite{ramanjotPlantDiseaseDetection2023}. Pendekatan ini memungkinkan deteksi penyakit secara otomatis tanpa ketergantungan penuh pada tenaga ahli maupun prosedur laboratorium yang kompleks.

\textcite{upadhyayDeepLearningComputer2025} memaparkan berbagai teknik \textit{deep learning} telah diterapkan untuk identifikasi penyakit tanaman. Teknik-teknik ini meliputi \textit{image recognition} (pengenalan citra) untuk klasifikasi jenis penyakit, \textit{object detection} (deteksi objek) untuk mendeteksi keberadaan penyakit, hingga \textit{vision-language models} (VLMs) yang mengintegrasikan pemrosesan informasi visual dan tekstual. Arsitektur VLM umumnya terdiri dari dua komponen utama, yaitu \textit{vision encoder} untuk mengekstraksi fitur dari data citra dan \textit{language model} untuk memproses informasi tekstual serta fitur visual yang telah diekstraksi \parencite{bordesIntroductionVisionLanguageModeling2024}. VLM sering dibangun di atas kemampuan \textit{large language models} (LLMs) \parencite{bordesIntroductionVisionLanguageModeling2024, quocVisionLanguageFoundationModel2025}, yaitu model AI yang dilatih pada data teks dalam skala masif untuk memahami, meringkas, dan menghasilkan bahasa. Sinergi ini memungkinkan sistem memproses data multimodal, yaitu data yang menggabungkan berbagai modalitas informasi seperti citra dan teks secara bersamaan. Identifikasi penyakit tidak hanya terbatas pada label klasifikasi visual, tetapi juga dapat disertai konteks deskriptif yang komprehensif.

Penerapan teknologi ini didorong oleh paradigma \textit{foundation models}, yaitu model yang dilatih pada skala data sangat masif dan beragam agar dapat diadaptasi untuk berbagai tugas hilir (\textit{downstream tasks}) \parencite{bommasaniOpportunitiesRisksFoundation2022}. Model-model ini, seperti GPT-3 dari \textcite{brownLanguageModelsAre2020} untuk pemrosesan teks atau CLIP dari \textcite{radfordLearningTransferableVisual2021} untuk tugas \textit{vision-language}, dirancang sebagai dasar yang dapat disesuaikan untuk berbagai aplikasi spesifik \parencite{radfordLearningTransferableVisual2021, liuSwinTransformerHierarchical2021}. Berkat kemunculan \textit{foundation models}, LLM telah diadopsi secara luas di berbagai domain. Kini prinsip tersebut dapat diperluas melalui VLM untuk menangani tantangan identifikasi penyakit tanaman sebagai perkembangan dari metode sebelumnya.

Tantangan utama pada identifikasi penyakit tanaman apabila melalui \textit{image recognition} terletak pada keterbatasan model konvensional. Model berbasis \textit{convolutional neural network} (CNN) masih lazim digunakan berkat kemampuannya mengekstraksi fitur secara efektif \parencite{gulmezComprehensiveReviewConvolutional2025}. Akan tetapi, CNN umumnya dilatih melalui \textit{supervised learning} untuk memetakan gambar ke kelas-kelas penyakit yang telah ditentukan sebelumnya \parencite{pacalSystematicReviewDeep2024}. Pendekatan ini menghasilkan klasifikasi yang kaku karena terbatas pada label yang tersedia selama pelatihan. Selain itu, model CNN memerlukan data latih dalam jumlah besar agar mampu mencapai performa tinggi \parencite{liuCrossModalDataFusion2025}. Pengumpulan citra penyakit tanaman dari berbagai jenis dengan kualitas yang memadai sulit dilakukan serta membutuhkan biaya dan waktu yang signifikan. Ketika jenis penyakit baru perlu ditambahkan, model harus dilatih ulang dengan data tambahan, yang berimplikasi pada kebutuhan komputasi tinggi.

Di lain sisi, pada \textit{object detection}, penelitian sebelumnya memperlihatkan bahwa YOLO (\textit{You Only Look Once}) masih mendominasi sebagai arsitektur utama untuk deteksi penyakit tanaman \parencite{ranaDetectionPlantDiseases2025}. Model ini mengadopsi pendekatan regresi tunggal yang memprediksi \textit{bounding box} dan probabilitas kelas secara langsung dari seluruh citra dalam satu langkah, sehingga sangat efisien untuk aplikasi \textit{real-time} \parencite[Redmon dkk., 2016 dalam][]{ranaDetectionPlantDiseases2025}. Dalam konteks identifikasi penyakit tanaman, YOLO telah mengalami perkembangan signifikan. YOLOv11 sebagai versi terbaru menawarkan peningkatan akurasi dan kecepatan inferensi dibandingkan pendahulunya \parencite{ranaDetectionPlantDiseases2025}. Analisis komprehensif \textcite{ramosComprehensiveAnalysisYOLO2025} membandingkan berbagai arsitektur YOLO, termasuk YOLOv8, YOLOv9, YOLOv10, YOLOv11, dan YOLOv12, untuk identifikasi penyakit daun tomat. Hasil analisis menunjukkan YOLOv11 memiliki performa paling efektif, dengan varian YOLOv11x mencapai mAP@50 sebesar 93,6\% dan mAP@50:95 sebesar 79,0\%. Dalam aplikasinya, YOLO umum diimplementasikan dalam skema \textit{two-stage detection} untuk penyakit tanaman \parencite{abhiramNovelTwoStageDeep2024a, yaoTwoStageDetectionAlgorithm2024}. Tahap pertama mendeteksi lokasi gejala pada daun dan tahap kedua mengklasifikasikan jenis penyakitnya. Pendekatan ini memungkinkan deteksi yang lebih presisi dengan kecepatan tinggi. Meskipun unggul dalam kecepatan, YOLO masih menghadapi tantangan serupa dengan CNN konvensional, yaitu ketergantungan pada data latih berlabel dalam jumlah besar dan kesulitan dalam mendeteksi penyakit yang tidak termasuk dalam kelas pelatihan \parencite{pacalSystematicReviewDeep2024}.

Pendekatan VLMs muncul sebagai solusi alternatif yang menjanjikan untuk mengatasi keterbatasan metode konvensional. Pada \textit{image recognition}, \textcite{radfordLearningTransferableVisual2021} mengembangkan CLIP (\textit{Contrastive Language-Image Pre-training}) yang memperkenalkan kemampuan \textit{zero-shot learning} dengan menggabungkan representasi visual dan teks melalui \textit{pre-training} pada skala data masif. \textit{Zero-shot learning} merupakan kemampuan model untuk mengenali kelas baru tanpa melihat contoh dari kelas tersebut selama pelatihan. Melalui \textit{zero-shot learning}, CLIP dapat mengklasifikasikan gambar tanpa memerlukan \textit{fine-tuning} pada kelas tertentu. Dalam domain penyakit tanaman, \textcite{quocVisionLanguageFoundationModel2025} mengembangkan SCOLD (\textit{Soft-target COntrastive learning for Leaf Disease identification}), model \textit{foundation vision-language} yang terinspirasi arsitektur CLIP. SCOLD dikembangkan bersama dataset LeafNet berisi pasangan gambar-teks yang secara khusus dirancang untuk identifikasi penyakit tanaman secara multimodal. Selain itu, \textcite{weiBenchmarkingInthewildMultimodal2024} juga mengembangkan dataset PlantWild berisi pasangan gambar-teks dan \textit{baseline} adopsi model CLIP untuk identifikasi penyakit tanaman. Dari dataset PlantWild, \textcite{weiSnapDiagnoseAdvanced2024} mengembangkan Snap and Diagnose untuk membangun sistem \textit{image retrieval} yang mendukung \textit{input} gambar dan teks, sehingga memungkinkan identifikasi penyakit melalui pencarian gambar serupa secara multimodal.

Selain itu pada \textit{object detection}, \textcite{pengLeafDiseaseImage2022} menunjukkan kombinasi YOLO untuk deteksi objek daun dengan \textit{deep metric learning} untuk \textit{image retrieval} mencapai akurasi 97,84\% pada \textit{closed-set identification} dan 89,09\% pada \textit{open-set identification}. Peng dkk. menekankan bahwa pendekatan ini dapat menggabungkan keunggulan YOLO dalam deteksi \textit{real-time} dengan kemampuan \textit{image retrieval} untuk mengatasi tantangan penambahan kelas baru. Sistem hanya perlu menambahkan gambar baru ke \textit{database retrieval} tanpa melatih ulang model, berbeda dengan metode klasifikasi konvensional yang memerlukan pelatihan ulang.

Pada perkembangan terbaru, transformasi signifikan terjadi dengan munculnya model \textit{vision-language} untuk \textit{object detection}. \textcite{liuGroundingDINOMarrying2025} mengembangkan Grounding DINO dengan menggabungkan detektor DINO berbasis Transformer dengan \textit{pre-training} yang memungkinkan deteksi berbasis \textit{prompt} teks bahasa natural. Grounding DINO membuka kemampuan \textit{open-vocabulary detection}, yakni deteksi objek berdasarkan deskripsi teks fleksibel tanpa batas kelas yang ditentukan sebelumnya. Hal ini memungkinkan deteksi objek baru tanpa data latih tambahan. Penelitian \textcite{singhFarmerChatScalingAIPowered2024} dan \textcite{weiBenchmarkingInthewildMultimodal2024} menunjukkan adaptasi Grounding DINO untuk deteksi tanaman dapat mencapai performa unggul bahkan dengan data terbatas. Selain Grounding DINO, \textcite{mindererScalingOpenVocabularyObject2024} memperkenalkan OWLv2 sebagai perkembangan dari OWL-ViT dengan teknik \textit{self-training} OWL-ST. OWLv2 menunjukkan performa superior dalam \textit{open-vocabulary detection}, di mana varian G/14 mencapai 47,2\% AP pada kelas \textit{rare} dataset LVIS, melampaui model sebelumnya secara signifikan. Lebih lanjut, evaluasi oleh \textcite{wuEvaluatingPerformanceOpenVocabulary2026} menunjukkan bahwa OWLv2 memiliki robustitas yang lebih baik dibandingkan Grounding DINO dan Detic dalam menghadapi degradasi kualitas citra, namun penggunaannya untuk deteksi penyakit tanaman masih belum banyak dieksplorasi.

Setelah berbagai pendekatan berbasis VLM dikembangkan, perkembangan terbaru menunjukkan upaya integrasi model deteksi objek berbasis \textit{deep learning} dengan LLM untuk identifikasi penyakit tanaman. Pada implementasi paling sederhana, \textcite{sbDesignImplementationAIPowered2024} mengembangkan sistem untuk deteksi penyakit tanaman sekaligus memberikan rekomendasi penanganan. Sistem ini memanfaatkan YOLOv5 untuk mengidentifikasi penyakit dan diintegrasikan dengan LLM dari Gemini API yang digunakan secara langsung untuk menghasilkan respons percakapan tentang penyakit tanaman. Akan tetapi, LLM digunakan secara langsung tanpa adanya augmentasi tambahan yang dapat membawa keterbatasan. \textcite{minaeeLargeLanguageModels2025} menekankan bahwa adopsi LLM pada domain tertentu memerlukan augmentasi agar dapat menghindari risiko halusinasi dan keluaran yang tidak konsisten. Salah satu pendekatan yang paling fleksibel untuk menghindari masalah ini adalah dengan memberikan sistem akses \textit{retrieval} (pencarian) ke data eksternal.

Oleh karena itu, pendekatan integratif ini selanjutnya dikembangkan oleh \textcite{mondalVisionMeetsLanguage2025} dengan menggabungkan deteksi objek dari YOLOv8 untuk identifikasi penyakit tanaman kopi secara cepat serta LLM yang diberi akses pengambilan ke \textit{knowledge base} (basis pengetahuan) tentang penyakit tanaman kopi. Akses tersebut diimplementasikan melalui kerangka \textit{retrieval-augmented generation} (RAG) sehingga respon LLM dilandasi dokumen rujukan yang relevan dan mampu menyajikan rekomendasi penanganan. Kendati demikian, pendekatan tersebut masih menyisakan keterbatasan berupa alur dari deteksi YOLOv8 ke bagian pengambilan dan generasi jawaban ditetapkan secara tetap, sehingga perilaku sistem relatif kaku dan kurang adaptif.

Menurut \textcite{anthropicBuildingEffectiveAgents2024}, sistem yang memanfaatkan LLM dengan pengambilan atau kapabilitas tambahan lain yang alurnya sudah ditentukan sejak awal dapat dikategorikan sebagai \textit{agentic system} (sistem agentik) berbasis \textit{workflow} (alur kerja). Kapabilitas tambahan tersebut disebut sebagai \textit{external tools} (alat eksternal) yang diintegrasikan untuk membantu LLM menyelesaikan tugas tertentu. Alternatif yang lebih adaptif dibandingkan \textit{workflow} adalah LLM \textit{agent} (agen), yakni sistem yang memungkinkan LLM memilih dan menggunakan alat secara fleksibel sesuai konteks tugas. Salah satu implementasi ini adalah pola \textit{reasoning and acting} (ReAct) yang diperkenalkan oleh \textcite{yaoReActSynergizingReasoning2023}. Agen LLM yang diimplementasikan dengan pola ReAct memungkinkan LLM melakukan penalaran terlebih dahulu (\textit{reasoning}) sebelum memutuskan tindakan (\textit{acting}) dengan memanggil dan menggunakan alat eksternal. ReAct telah menjadi salah satu bentuk implementasi agen paling populer \parencite{yuSurveyAgentWorkflow2025}, seiring dengan munculnya LLM modern yang mendukung penggunaan alat eksternal, seperti keluarga Gemini 2.5 \parencite{comaniciGemini25Pushing2025} dan Gemini 3 \parencite{google2025}. Dengan demikian, sistem agentik yang mengadopsi pola ReAct dapat menjadi solusi untuk mengembangkan sistem yang lebih adaptif dan fleksibel.

Akan tetapi, peningkatan fleksibilitas pada sistem agentik ini turut membawa kompleksitas baru dalam proses validasi. Evaluasi kinerja sistem agentik berbasis LLM menghadirkan tantangan tersendiri karena luarannya berupa bahasa alami yang kompleks, cair, dan tidak terstruktur \parencite{liuGEvalNLGEvaluation2023}. Metrik evaluasi heuristik BLEU (\textit{Bilingual Evaluation Understudy}) dan ROUGE (\textit{Recall-Oriented Understudy for Gisting Evaluation}) kurang bisa menangkap nuansa semantik dan koherensi dari jawaban yang dihasilkan \parencite{saiSurveyEvaluationMetrics2022}. Penilaian oleh manusia sejatinya merupakan standar emas untuk menangkap nuansa tersebut, namun proses ini cenderung lambat, mahal, dan sulit dilakukan dalam skala besar. Oleh karena itu, pendekatan \textit{LLM-as-a-judge} diperkenalkan sebagai solusi alternatif untuk memanfaatkan kemampuan LLM yang kuat dalam menilai kualitas luaran model lain menyerupai penilaian manusia \parencite{liGenerationJudgmentOpportunities2025}. Akan tetapi, penggunaan LLM sebagai evaluator juga tidak luput dari kelemahan, seperti adanya bias posisi (\textit{positional bias}) di mana model cenderung memfavoritkan jawaban yang muncul pada urutan pertama \parencite{wangLargeLanguageModels2024}. Untuk mengatasi keterbatasan tersebut dan meningkatkan akurasi evaluasi, berbagai kerangka kerja (\textit{framework}) telah dikembangkan. \textcite{liuGEvalNLGEvaluation2023} memperkenalkan G-Eval yang memanfaatkan teknik \textit{chain-of-thought} (CoT) untuk penilaian yang lebih selaras dengan preferensi manusia. Selain itu, \textit{framework} seperti RAGAS \parencite{esRagasAutomatedEvaluation2025} dan DeepEval \parencite{ipDeepeval2026} kini tersedia untuk mengevaluasi komponen sistem agentik dengan lebih andal dan terukur.

Berdasarkan tantangan dan peluang yang telah diuraikan, penelitian ini mengusulkan pengembangan sistem agentik untuk identifikasi penyakit tanaman. Sistem ini akan mengintegrasikan VLM sebagai alat pada sistem untuk mengatasi keterbatasan pendekatan konvensional. Secara spesifik, penelitian ini akan mengadopsi model \textit{foundation} seperti SCOLD untuk identifikasi penyakit berbasis \textit{image-text retrieval} atau pencarian gambar dan teks, serta mengombinasikan YOLO dengan OWLv2 untuk membantu deteksi objek tanaman dalam identifikasi penyakit. Lebih lanjut, mekanisme \textit{retrieval-augmented generation} (RAG) akan dimanfaatkan untuk memberikan akses ke basis pengetahuan yang relevan. Keseluruhan kapabilitas ini akan diintegrasikan ke dalam sebuah agen LLM yang dirancang dengan pola \textit{reasoning and acting} (ReAct), memungkinkannya untuk memilih dan menggunakan alat secara adaptif dalam proses identifikasi hingga pemberian rekomendasi penanganan. Selain itu, sistem ini akan dilengkapi dengan antarmuka berbasis web untuk memfasilitasi interaksi pengguna. Untuk mengevaluasi performa sistem agentik yang dikembangkan, penelitian ini akan menggunakan pendekatan \textit{LLM-as-a-judge} dengan kriteria evaluasi yang relevan. Melalui integrasi ini, penelitian ini diharapkan dapat menghasilkan sistem identifikasi penyakit tanaman yang tidak hanya akurat dan komprehensif, tetapi juga fleksibel dan adaptif dalam menangani berbagai kasus identifikasi penyakit tanaman.

\section{Rumusan Masalah Penelitian}
Berdasarkan uraian latar belakang, maka permasalahan yang akan dibahas dalam penelitian ini dapat dirumuskan sebagai berikut: 
\begin{enumerate}[nosep]
    \item Bagaimana pengembangan sistem agentik yang mengintegrasikan deteksi multimodal, mekanisme \textit{Retrieval-Augmented Generation} (RAG), dan pola \textit{Reasoning and Acting} (ReAct) untuk identifikasi penyakit tanaman?
    \item Bagaimana perancangan dan pengembangan antarmuka pengguna berbasis web dilakukan untuk memfasilitasi interaksi dengan sistem agentik?
    \item Bagaimana evaluasi performa sistem agentik dilakukan melalui pendekatan \textit{LLM-as-a-judge}?
\end{enumerate}

\section{Tujuan Penelitian}
Tujuan dari penelitian ini adalah: 
\begin{enumerate}[nosep]
    \item Mengembangkan sistem agentik yang mengintegrasikan deteksi multimodal, mekanisme RAG, dan pola ReAct untuk identifikasi penyakit tanaman.
    \item Merancang dan mengembangkan antarmuka pengguna berbasis web untuk memfasilitasi interaksi dengan sistem agentik.
    \item Mengevaluasi performa sistem agentik yang dikembangkan dengan pendekatan \textit{LLM-as-a-judge}.
\end{enumerate}

\section{Ruang Lingkup Penelitian}
Ruang lingkup yang menjadi batasan penelitian ini adalah: 
\begin{enumerate}[nosep]
    \item Karakteristik penyakit tanaman yang diidentifikasi oleh sistem mengikuti cakupan dan variasi gejala yang tersedia dalam dataset penyakit tanaman yang digunakan sebagai rujukan penelitian.
    \item Sistem yang dikembangkan tidak sepenuhnya mendukung bahasa Indonesia karena beberapa proses internal dan keluarannya dioptimalkan untuk bahasa Inggris menyesuaikan dataset yang digunakan. Akan tetapi, keluaran akhir sistem tetap dapat disajikan dalam bahasa Indonesia.
    \item Evaluasi sistem dibatasi pada metode \textit{offline evaluation} menggunakan dataset terkurasi untuk \textit{benchmarking} dan deteksi regresi, sedangkan evaluasi \textit{online} melalui interaksi pengguna secara \textit{real-time} belum dilakukan.
    \item Evaluasi lapangan secara langsung dengan data \textit{real-time} atau validasi jangka panjang pada praktik pertanian belum dilakukan.
    \item Luaran penelitian berupa sistem prototipe sebagai bukti konsep (\textit{proof of concept}) dan bukan merupakan aplikasi yang siap untuk komersialisasi.
\end{enumerate}

\section{Manfaat Penelitian} 
Penelitian ini diharapkan dapat memberikan kontribusi, baik dari segi pengembangan ilmu pengetahuan maupun penerapan praktis di lapangan.

\subsection{Manfaat Teoritis}
Penelitian ini berkontribusi pada pengembangan ilmu pengetahuan terkait integrasi model deteksi objek multimodal dan sistem agentik. Kajian ini memperkaya literatur mengenai penerapan pola ReAct serta RAG di bidang pertanian. Selain itu, metodologi \textit{LLM-as-a-judge} yang diterapkan dapat menjadi referensi validasi sistem agentik dengan luaran bahasa alami.

\subsection{Manfaat Praktis}
Selain manfaat teoritis, penelitian ini diharapkan memberikan manfaat praktis, antara lain:
\begin{enumerate}[nosep]
    \item Bagi petani, sistem yang dikembangkan diharapkan dapat dimanfaatkan sebagai asisten cerdas untuk mendeteksi penyakit tanaman secara dini dan akurat. Informasi diagnosis serta rekomendasi penanganan yang diberikan berpotensi membantu pengambilan keputusan yang lebih cepat tanpa ketergantungan penuh pada tenaga ahli.
    \item Bagi masyarakat umum, teknologi ini diharapkan dapat digunakan untuk memperoleh wawasan mengenai kondisi kesehatan tanaman. Hal ini berpotensi mendukung kemandirian masyarakat dalam merawat tanaman di lingkungan sekitar serta meningkatkan kesadaran akan pentingnya pengendalian penyakit tanaman.
    \item Bagi peneliti dan pengembang sistem, arsitektur sistem agentik yang diusulkan dapat menjadi rujukan bagi pengembangan solusi teknologi serupa di masa depan. Pendekatan integrasi berbagai model kecerdasan buatan serta kerangka evaluasi yang diterapkan dapat diadopsi dan disempurnakan untuk berbagai kasus penggunaan lain.
\end{enumerate}



